%!TEX root = ../../main.tex
% ============================================================
% 力学作图 / 弹簧测力计
% 本文件由 main.tex 通过 \input 引入，请勿单独编译
% 重构日期: 2026-04-11
% ============================================================

\subsection{解答：探究弹簧测力计原理及刻度板绘制}

\vspace{0.6cm}

\textbf{（1）从表中的数据可以看出弹簧的伸长量 $\Delta l$ 与拉力 $F$ 的关系为：}\\
\underline{\quad 在弹性限度内，弹簧的伸长量与受到的拉力成正比。\quad}

\vspace{0.4cm}

\textbf{（2）小明继续实验，得到如下数据：} \\
从这次数据看出：拉力达到 \underline{\quad $\textbf{11}$ \quad} N 时，弹簧的伸长量和拉力的关系就改变了。因此，该弹簧测力计的最大测量值为 \underline{\quad $\textbf{10}$ \quad} N。

\vspace{0.4cm}

\textbf{（3）请你根据表中的数据，在图中用刻度尺画出这个弹簧测力计的刻度板（分度值为 1N）。}

\vspace{0.4cm}

\begin{center}
\begin{tikzpicture}[>=Stealth, line join=round, line cap=round, scale=1.0]
  % ---------------------------------------------------
  % 1. 测力计面板外轮廓 (含上下半圆凸起，且加宽边距防数字重叠)
  % ---------------------------------------------------
  \draw[black, line width=1.0pt, fill=white] 
    (-1.0, 0.8) -- (-0.35, 0.8) arc(180:0:0.35) -- (1.0, 0.8) -- 
    (1.0, -4.6) -- (0.35, -4.6) arc(360:180:0.35) -- (-1.0, -4.6) -- cycle;
    
  % ---------------------------------------------------
  % 2. 中间的直槽 (上下两端圆润过渡)
  % ---------------------------------------------------
  \draw[black, line width=0.8pt, fill=gray!5] 
    (-0.15, -4.3) -- (-0.15, 0.3) arc(180:0:0.15) -- 
    (0.15, -4.3) arc(360:180:0.15) -- cycle;
    
  % ---------------------------------------------------
  % 3. 顶部的单位 'N' 
  % ---------------------------------------------------
  \node[font=\large] at (0, 0.9) {N}; 
  
  % ---------------------------------------------------
  % 4. 绘制 0 刻度线 (直接从右侧槽边伸出，高度契合原图外壳)
  % ---------------------------------------------------
  \draw[black, line width=1.0pt] (0.15, 0) -- (0.4, 0); 
  \node[right, font=\small, inner sep=2pt] at (0.4, 0) {0};

  % ---------------------------------------------------
  % 5. 严格按物理极距 (0.4cm) 补充的蓝色答题层刻度
  % ---------------------------------------------------
  \foreach \F in {1, ..., 10} {
    \pgfmathsetmacro{\y}{-0.4 * \F}
    % 画分度值为 1N 的主刻度线
    \draw[blue!80!black, line width=1.0pt] (0.15, \y) -- (0.4, \y);
    \node[right, text=blue!80!black, font=\small, inner sep=2pt] at (0.4, \y) {\F};
  }
\end{tikzpicture}
\end{center}

\vspace{0.6cm}
\hrule
\vspace{0.4cm}

{\small\color{gray}
\textbf{解题分析与说明：}\\
1. \textbf{定性归纳第一题：}在 $1\sim 7\,\text{N}$ 的区间内，拉力 $F$ 每增加 $1\,\text{N}$，伸长量 $\Delta l$ 恒定增加 $0.40\,\text{cm}$。这验证了胡克定律的经典表述：“在弹性限度内，弹簧的伸长量与受到的拉力成正比”。\\
2. \textbf{寻找拐点破第二题：}拉力到 $8, 9, 10\,\text{N}$ 时，其伸长量 $3.20, 3.60, 4.00\,\text{cm}$ 依然吻合正比例。但当拉力达到 \textbf{$11\,\text{N}$} 时，伸长量实测值变为 $4.50\,\text{cm}$（理论应为 $4.40\,\text{cm}$），说明正比例关系被彻底打破。因此，测量上限必须卡死在正比区间的最高边缘也就是 \textbf{$10\,\text{N}$}。\\
3. \textbf{严谨绘制标线：}明确了最大量程为 $10\,\text{N}$，且分度值为 $1\,\text{N}$，就意味着需要等比例精确无误地画出 $10$ 根刻度线。每根线的间隔长度必须严格依据实测物理极距 $0.40\,\text{cm}$。
}

\vspace{0.6cm}

{\small\color{blue!70!black}
\textbf{代码自动化处理解析：}\\
\textbf{1. 物理厘米级等距映射刻度：}不同于手绘粗略示意，代码可以直接绑定物理真实步距。在默认 $1.0 = 1.0\,\text{cm}$ 的比例系中，循环逻辑 $\texttt{\textbackslash pgfmathsetmacro\{\textbackslash y\}\{-0.4 * \textbackslash F\}}$ 驱动画笔，确保在纸面上呈现出的真实物理 $0.4\,\text{cm}$ 间隔极距跨步，还原了实验中游标下沉的物理真实距离。\\
\textbf{2. 循环量程控制锁：}刻度条数的决定并非随意排布，而是由题干第二问的“弹性限度安全区域”强力关联绑定。\texttt{\textbackslash foreach} 循环被强制封死在最大值 \texttt{10} 处，阻断了刻度越界超画的安全隐患。\\
\textbf{3. 答题区色彩直观防线：}为了突显被学生后期补绘卷面充填出来的内容，我专门动用了 \texttt{blue!80!black} 为所有新生刻度线和标识数字进行深度强着色。这不仅展现了答题前后的步骤层叠感，更使正确刻划痕变得格外醒目。
}

%% ==============================
%% 附加图：m-V (质量-体积) 图像基础网格
%% ==============================
